\scanpage{47}
\mspage{- 15 -}
n\textsuperscript{o} 276

linéaire des $\alpha_{rs}X^r$ si $r\leq n$, et des $\alpha_{ns}X^r$ si $r>n$, d'où évidemment notre assertion.

(La démonstration précédente ne semble pas présenter d'avantage très décisif par rapport à la démonstration classique \ldots)

\uline{Corollaire} 1 - \uline{Soit} $K$ \uline{un corps} ; \uline{l'anneau de polynômes} $K[X_1,..,X_r]$ \uline{est noethérien}.

On raisonne par récurrence sur $r$.

\margin{\handwritten{?}}\uline{Corollaire} 2 - \uline{Soit} $A$ \mbox{\settowidth{\unitlength}{un anneau commutatif noethérien,}\handwritten{\tikz[overlay,baseline,x=\unitlength,y=1ex]{\draw[thin] (.18,1.55) .. controls (-.1,1.2) and (-.025,-.68) .. (.16,-.62) .. controls (.52,-.37) and (.64,-1.42) .. (.86,-.99) .. controls (1.15,-.47) and (1.09,1.60) .. (.76,1.75);}}\uline{un anneau commutatif noethérien},} \uline{L'anneau de séries formelles} $A[[X]]$ \uline{est noethérien}.

Filtrons $A[[X]]$ par les puissances de l'idéal $(X)$ ; on obtient un anneau séparé et complet (Corollaire de la Proposition 5) ; or l'anneau gradué associé, isomorphe à $A[X]$, est noethérien ; donc $A[[X]]$ est noethérien en vertu du corollaire 2 de la proposition 7.

\uline{Corollaire} 3 - \uline{Soit} \mbox{\settowidth{\unitlength}{$K$ un corps ;}\handwritten{\tikz[overlay,baseline,x=\unitlength,y=1ex]{\draw[thin] (0,.47) .. controls (-.10,1.75) and (.79,2.25) .. (1.01,1.36) .. controls (1.23,.2) and (.99,-.93) .. (.64,-1.01) .. controls (.26,-1.35) and (-.08,-.89) .. (0,.47);}}$K$ \uline{un corps} ;} \uline{l'anneau de séries formelles} $K[[X_1,..,X_r]]$ \uline{est noethérien}.
% margin note written sideways
\margin{\handwritten{\tiny\setlength{\arraycolsep}{.2pt}
$\begin{array}{ccc}
\overbrace{(\underbrace{0,1,\infty})}&
\tikz[baseline=-.5ex,scale=.15]{\draw (0,-1.3)--(0,1.3) (1.2,1.15)--(-1.2,-1.35);}&
\overbrace{(2,\underbrace{1/2,-1})}\\
&(x',x''\unsure{,})&
\end{array}$}}

Cela résulte du corollaire précédent par récurrence sur $r$.

Le corollaire 2 résulte aussi du suivant :

\uline{Corollaire} 4 - \uline{Soient} $A$ \uline{un anneau commutatif et} $m$ \uline{un idéal de} $A$. \uline{Supposons} $A/m$ \uline{noethérien}, $m$ \uline{à engendrement fini}, \uline{et} $A$ \uline{séparé et complet pour la topologie} $m$\uline{-adique}. \uline{Alors} $A$ \uline{est noethérien}.

En effet, si $m$ est engendré par des éléments $x_i$ en nombre fini dont nous noterons $\bar x_i\in G_1(A)$ les formes initiales, il est clair que tout élément de $G(A)$ est un polynôme en les $\bar x_i$, à coefficients dans l'anneau $G_0(A)=A/m$ ; $G(A)$ est donc noethérien, d'ou le résultat.

\uline{Corollaire} 5 - \uline{Soit} $g$ \uline{une algèbre de Lie sur un anneau commutatif noethérien} $K$. \uline{L'algèbre enveloppante de} $g$ \uline{est un anneau noethérien} (\uline{à gauche et à droite})
% archive sheet number: 46
